\section{结论}
\label{sec:conclusion}

本文提出一个关于 CLIP 的观察：异常相关的线索其实已经存在于 CLIP 的高频分量中，缺的不是能力，而是一个逐图的正常参照。据此我们提出一种训练无关的方法——读出高频、用这个 CLIP 之外的信号选择证据、逐图估计正常参照——把一个直接使用会崩盘的线索变成可靠的零样本定位。

\paragraph{局限。}
相对 CLIP-only 原型自适应（SelfRef），总均值增益较小；本文的机制主要由崩盘负控、参照的逐图性与盲区召回三条独立证据支撑，而非依赖对 SelfRef 的均值碾压。此外，MPDD、BTAD 与 DTD 使用了逐数据集设置，逻辑异常仍然困难。

\paragraph{未来工作。}
将“逐图正常参照”推广到其他冻结编码器；以及与轻量适配相结合，在保持训练无关优势的同时进一步提升判别力。
